\section{The actual tuple link to the native ordinary determinant ring}
\label{sec:cubic-tuple-native-link}

The original \path{TwLink.chib} requires powers of one character,
and \path{TwLink.ker} uses the kernel of that character.  The cubic
tuple has a different joint kernel.  This section constructs the
actual quotient map and checks every field of the native
\path{Pair} for that tuple.  No assertion that the old
\path{TwLink} record has been instantiated is made.

\paragraph{Actual arithmetic data.}
Put \(p=269\).  Let \(H/F\), \(F=H^+\), be one of the
actual CM hosts of AC-G93-fill, at a fixed finite cyclotomic
layer.  Retain its proved finite-unramified CM extension,
the splitting over \(F\) of every coefficient-prime place and
every prescribed bad place, and its coefficient ring
\(\mathcal O\), a finite extension of \(\mathbb Z_p\) with
residue field \(k\).  Fix a common algebraic closure.
Write \(\mathcal G_F,\mathcal G_H\) for the full absolute
Galois groups, \(c\) for the chosen complex conjugation, and
\(\eta:\mathcal G_F\to\{\pm1\}\) for the CM character.

On \(\mathcal G_H\) the residual tuple is
\[
 \boldsymbol\theta=(\bar\chi^j\gamma,1,
                         \bar\chi^j\gamma^{-1}),\qquad j=1\text{ or }2.
 \tag{G97.1}
\]
It is the actual retained tuple, with the local separation
used in AC-G94.  Let \(N\) be the closed normal subgroup
of \(\mathcal G_F\) generated by
\[
 I_w\cap\bigcap_i\ker\theta_i
       \quad(w\text{ a finite place of }H,\ w\nmid p).
 \tag{G97.2}
\]
This family includes all conjugate places.  It is stable under
\(c\), since conjugation reverses and inverts the residual
characters.  The group lies in \(\mathcal G_H\).
Put \(\Gamma=\mathcal G_F/N\) and
\(G=\mathcal G_H/N=\ker(\Gamma\to\{\pm1\})\).
This is the actual AC-G90/G96 away-kernel quotient.
The residual characters and the cyclotomic character descend
to these groups.

Choose the already constructed native \path{HidaSetting} on
\(H\), with auxiliary place \(v_0\), and its finite set
\(S=s.S\) of places of \(F\).  It contains the places above
\(p\), every place where (G97.1) ramifies, every retained
automorphic bad place, and \(v_0\).  The auxiliary place is
outside the old ramification.  Retain a native \path{Lev}
with \(V.S'=S\), as in AC-G89 and AC-G93-fill.
Thus its native group is
\[
 \mathcal G_S=\mathcal G_F/N_S,
 \qquad N_S=\overline{\langle\!\langle I_v:v\notin S\rangle\!\rangle},
 \qquad \mathcal G_{S,0}=\ker\eta_S .
 \tag{G97.3}
\]
The local groups in the native \path{GFrame} are the actual
decomposition groups over \(H\); they are also the local
groups used in AC-G94.  Their maps into \(G\) are the
quotient maps.  This specifies the local comparison, rather
than replacing a full local group by its image.

\begin{proposition}[AC-G97-link: the actual joint-kernel tuple link]
\label{AC-G97-link}
There is a continuous surjection
\[
 \pi:\mathcal G_{S,0}\longrightarrow G
 \tag{G97.4}
\]
compatible with complex conjugation, the cyclotomic character
and every chosen local map.  Its kernel is generated by the
images of the away-\(p\) joint residual-kernel inertia in
(G97.2).  In particular this is a proved joint-kernel
replacement of the original single-character kernel clause.
Equivalently, after the native residual data below have been
defined, \(\ker\pi\) is the closed normal, conjugation-stable
subgroup generated by the native groups \(\operatorname{Ip}(u)\)
and \(I_0\).

Put \(\kappa=\epsilon^{-1}:\mathcal G_{S,0}\to\mathcal O^\times\),
using the native cyclotomic character.  Then the characters
\[
 \bar t_i=(\theta_i\circ\pi)\bar\kappa
 \tag{G97.5}
\]
are the actual ordered native residual tuple.  They define
a native \path{Resid}, with local characters the restrictions
of (G97.5) and with \path{th} equal to one.
For every coefficient-prime local group the literal native
subgroup \path{mu} equals the closed commutator subgroup.
For every prescribed away inertia group, the image under
\(\pi\) is finite of order prime to \(p\), and \(\kappa\)
is trivial.  The same assertions give triviality of both
\(\pi\) and \(\kappa\) on the auxiliary inertia \path{I0}.
\end{proposition}

\begin{proof}
At a finite place outside \(S\), the CM extension is
unramified, and all residual characters are unramified.
Full inertia therefore belongs to (G97.2).  Hence
\(N_S\subset N\).  The universal property of the profinite
quotient gives a continuous surjection
\(\mathcal G_S\to\Gamma\).  Both quotient CM characters
come from \(\eta\), so restriction to their kernels is
surjective: a lift of an element with CM value one already
has CM value one.  This gives (G97.4).
Its kernel is \(N/N_S\), with exactly the stated closed
generation.  All conjugation and local identities follow
from the same quotient map on the absolute group.

The cyclotomic character is unramified at every finite place
away from \(p\), so it kills \(N\).  The native character
on \(\mathcal G_S\) is the same descended character, by
the defining \path{Lev} cyclotomic compatibility.
Thus \(\pi\) preserves it.
On \(\mathcal G_{S,0}\), conjugation fixes this character,
and consequently
\[
 \kappa(g)\kappa(cgc^{-1})=\epsilon(g)^{-2}.
 \tag{G97.6}
\]
Every local comparison may be taken to be the identity on
the chosen decomposition group; its inertia subgroup is
unchanged.

For an away inertia group \(I_w\), the homomorphism to
\(G\) kills \(I_w\cap\bigcap_i\ker\theta_i\).
It therefore factors through the finite image of the tuple
of characters.  This image has order prime to \(p\).
Moreover \(\kappa|_{I_w}=1\).  At \(v_0\) the tuple is
unramified as well, so the entire inertia group is killed.
Passing to its image under the native quotient proves the
assertions about \path{Iu} and \path{I0}; no semisimplified
inertia statement is being substituted for actual killing.

We verify the coefficient-prime torsion clause explicitly.
In the retained hosts, each coefficient-prime completion
is obtained from the initial unramified quadratic
\(\mathbb Q_p\)-completion by unramified \(p\)-power
extensions, extensions of \(p\)-power degree from the
cyclotomic layers, and the selected-place extensions that
split completely at these primes.  Thus its degree over
\(\mathbb Q_p\) has no factor \(67\).
It cannot contain \(\mu_p\), since
\([\mathbb Q_p(\mu_p):\mathbb Q_p]=268=4\cdot67\).
This is an argument about the actual local fields, not the
degree of an arbitrary non-Galois global host.

Local class field theory identifies the topological
abelianization of the absolute local group with the
profinite completion of the multiplicative local field.
Its valuation factor \(\widehat{\mathbb Z}\) is torsion-free.
The prime-to-\(p\) residue-unit group has no \(p\)-torsion;
the \(p\)-primary torsion in the principal units consists
exactly of the \(p\)-power roots of unity.
The latter is zero here.  Thus the \(p\)-primary torsion
subgroup of the topological abelianization is zero.
The source definition of \path{GFrame.mu} is its preimage,
which is precisely the closed commutator subgroup.

The characters (G97.5) have finite image and open kernels.
They are distinct, since one retained local element has
three distinct \(\theta_i\)-values and the common scalar
\(\bar\kappa\) does not change that fact.
The tuple identity
\(\theta_i(cg^{-1}c^{-1})=\theta_{2-i}(g)\), together
with (G97.6), gives the required native residual
polarization: on each matrix, the native group-algebra
star becomes transpose conjugated by the antidiagonal
matrix \(J\).  Determinants are unchanged, also after
arbitrary coefficient base change.
The local residual characters are the restrictions of
(G97.5).  Set their prescribed \path{th} characters on
\path{mu} to one.  The \path{th_comm} field is then
literal.  This verifies all fields of \path{Resid}.
Indeed any existing native \path{Resid} with these global
and local characters has the same \path{th}: its
\path{th_comm} forces it to be one when \path{mu} is the
closed commutator.  Thus these are the same residual data
used for the actual native support construction.

We finally verify the finite native description of the
kernel.  Inertia outside \(S\) is already killed in
\(\mathcal G_S\), and no coefficient-prime inertia occurs
among the generators of \(N\).  At each remaining
prescribed place the extension \(H/F\) splits.  Its two
places are exchanged by complex conjugation, and any
choice of decomposition group differs by conjugation
inside \(\mathcal G_H\).  On their inertia \(\kappa=1\).
Consequently the image of
\(I_w\cap\bigcap_i\ker\theta_i\) is exactly
\(I_u\cap\bigcap_i\ker\bar t_i=\operatorname{Ip}(u)\):
an element of the latter has an inertia preimage, and
its residual character values are unchanged by that
preimage.  At \(v_0\) the whole inertia image is \(I_0\).
Taking closed normal, conjugation-stable generation
therefore gives precisely \(N/N_S\), not merely a
subgroup contained in it.  This proves the asserted
joint-kernel analogue of the native kernel clause.
\end{proof}

\paragraph{Rings and notation.}
Let \(R\) be the AC-G94-space ring for \(G\), and let
\(B\subset A\subset R\) be its actual AC-G95 algebras.
Write \(r:G\to\operatorname{GL}_3(A)\) for the restricted
universal representation and
\(\zeta_{v,i}:G_v\to B^\times\) for its local characters.
They are units in \(B\), because their inverses are the
values at inverse group elements.  Define
\[
 \widetilde r(g)=\kappa(g)r(\pi g),\qquad
 \widetilde\zeta_{v,i}(h)
      =\kappa(\phi_v h)\zeta_{v,i}(h),
 \tag{G97.7}
\]
where \(\phi_v\) is the native local map.  The first
representation has entries in \(A\), and the local
characters have values in \(B^\times\).

\begin{theorem}[AC-G97-pair: the native Pair and its exact map direction]
\label{AC-G97-pair}
The determinant of (G97.7) descends to a degree-three
native \path{GDet} \(D_B\) over \(B\).  The pair
\((D_B,\widetilde\zeta)\) is an actual native
\path{Pair} for (G97.5), satisfying every displayed field
of that definition.

If \(\mathscr B\) is the native universal ordinary
determinant ring for these residual data, there is a unique
\(\mathcal O\)-algebra homomorphism
\[
                  \Phi:\mathscr B\longrightarrow B
 \tag{G97.8}
\]
carrying its universal determinant and local characters
to \((D_B,\widetilde\zeta)\).  The map is local and
continuous.  No surjectivity or injectivity of \(\Phi\)
is asserted.

For a native \path{LamObj} \(\Lambda\), its map to this
pair is exactly the composite
\[
        \Lambda\longrightarrow\mathscr B
                   \xrightarrow{\Phi}B .
 \tag{G97.9}
\]
Finally, if a native ideal \(J_{\rm aut}\subset\mathscr B\)
is proper, put \(M=\Phi(J_{\rm aut})B\).  Then
\[
 M\subset\mathfrak m_B,\qquad
 MA\subset x_0,\qquad
 MA\subset\ker\xi
 \quad\hbox{for every AC-G95 tuple \(\mathrm{FPt}\) map }\xi.
 \tag{G97.10}
\]
In particular this is actual closed-invariant-fiber
containment for the AC-G95 polynomial paths.  It gives
no containment for a prescribed DVR curve whose
weight-zero characters vary.
\end{theorem}

\begin{proof}
\emph{1. Integral descent of the polynomial determinant law.}
For a finite list \(g_1,\ldots,g_m\), expand
\[
 \det\left(\sum_{a=1}^m T_a\widetilde r(g_a)\right)
                         \in A[T_1,\ldots,T_m].
 \tag{G97.11}
\]
Every term of its determinant expansion has total torus
weight \(\sum_i(w_i-w_{\sigma(i)})=0\).
The factors \(\kappa(g_a)\) are scalars from
\(\mathcal O\).  Hence every coefficient lies in
\(A_0=B\).  These coefficient polynomials define a
natural degree-three law on each finite support in the
group algebra over any \(B\)-algebra.  Enlarging the
support, changing coefficients, multiplication and
homogeneity satisfy the determinant identities:
write each identity as a polynomial in finitely many
independent variables and check it in \(A\), where it
is the ordinary matrix determinant identity.
The inclusion \(B\hookrightarrow A\) is injective,
so every coefficient identity already holds in \(B\).
This constructs the native \path{GDet} over \(B\).
It proves its naturality for arbitrary coefficient
base change without assuming that base change is flat.
After extending to \(A\), it is exactly
\(\operatorname{GDet.ofRep}(\widetilde r)\).

\emph{2. Continuity and residual data.}
The literal \path{GDet.IsCont} asks for adic continuity
of the characteristic-polynomial coefficients of every
group element.  They are polynomial expressions in
the continuous entries over \(R\), so are
\(\mathfrak m_R\)-adically continuous.  Their differences
have weight zero.  The AC-G95 filtration bound gives
\[
 \mathfrak m_R^{3m+2}\cap B\subset\mathfrak m_B^m.
 \tag{G97.12}
\]
Indeed AC-G95 first gives membership in
\(\mathfrak m_B^m A\), and its algebraic grading shows
\(\mathfrak m_B^m A\cap A_0=\mathfrak m_B^m\).
Thus the required continuity is in the actual
\(\mathfrak m_B\)-topology.  The same argument applies
to \(\zeta_{v,i}\); multiplication by the continuous
\(\mathcal O\)-valued \(\kappa\) preserves it.

The residue map of \(B\) is the restriction of the
residue map of \(R\).  Reducing (G97.11) there gives
the determinant of the diagonal tuple (G97.5), as a
polynomial law.  This is exactly \path{Pair.red}.
The local characters reduce to the native \path{psib}.
They kill the closed local commutator because they
are continuous characters into the Hausdorff group
\(B^\times\).  AC-G97-link identifies this subgroup
with \path{mu}.  Hence their values there are one,
exactly the prescribed \path{th}.  This checks all
three \path{IsPsi} fields: \path{cont}, \path{red},
and \path{mu}.

\emph{3. The native polarization field.}
The universal polarization and (G97.6) give
\[
 \widetilde r(g)^tJ\widetilde r(cgc^{-1})
                           =\epsilon(g)^{-2}J.
 \tag{G97.13}
\]
The source \path{GFrame.starOp} in rank three sends
\([g]\) to \(\epsilon(g)^{-2}[(cgc^{-1})^{-1}]\).
Under the matrix representation its image is therefore
\(J\widetilde r(g)^tJ\), by (G97.13).
Linearity gives the same transpose-conjugation formula
for every group-algebra element, over every coefficient
base change.  Its characteristic polynomial and
determinant are unchanged.  In particular every
\path{lam i} on the two elements in \path{Pair.P}
is equal.  The coefficientwise descent in Step 1
places this equality in \(B\) itself.

\emph{4. The ordinary and away-inertia fields.}
For the native ordinary product, the matrix of its
\(i\)-th factor is
\[
 \widetilde r(\phi_v h_i)
       -\widetilde\zeta_{v,i}(h_i)I
   =\kappa(\phi_v h_i)
       \bigl(r(\iota_v h_i)-\zeta_{v,i}(h_i)I\bigr).
 \tag{G97.14}
\]
Its ordered product is zero by AC-G94.  Multiplication
by \(\widetilde r(h)\) still gives the zero matrix;
its trace is zero.  This is precisely \path{Pair.Ord}.

The literal native inertia kernel is
\[
 \operatorname{Ip}(u)=I_u\cap\bigcap_i\ker\bar t_i.
 \tag{G97.15}
\]
For \(g\) in this group, choose an inertia preimage
in the absolute local group.  Its \(\kappa\)-value is
one, so (G97.5) and (G97.15) say that every
\(\theta_i\)-value is one.  The defining quotient
(G97.2) kills that element, proving
\(\pi(g)=1\) and \(\widetilde r(g)=I\).
Consequently the matrix of \([h]([g]-1)\) is zero
for every \(h\), and its trace vanishes in \(B\).
This is \path{Pair.I}.  It uses the full actual
joint inertia kernel, not a conclusion about a
semisimplified Weil--Deligne parameter.
All fields of the native \path{Pair} are now checked.

\emph{5. Universality and the exact support scope.}
The ring \(B\) is an \path{IsCO} test ring by AC-G95.
The defining native universal property therefore gives
the unique map (G97.8), with exactly its stated map
direction.  Its automatic locality and continuity
follow from the AC-G94 local-map argument, applied
to \(\mathscr B\) and \(B\).
If existence of \(\mathscr B\) has not already been
supplied by AC-G89, the retained source proposition
\path{ProModularRingExists}, applied to this actual
\path{Pair}, supplies it.  That proposition remains
an explicitly retained source input, not a theorem
derived here from the analytic hypothesis alone.

The source \path{LamObj.toPair} is defined by its
universal property for the restrictions to inertia.
Both maps in (G97.9) carry every universal inertial
character to the restrictions of
\(\widetilde\zeta_{v,i}\).  Uniqueness in that same
universal property proves (G97.9).

If \(J_{\rm aut}\) is proper, it is contained in
\(\mathfrak m_{\mathscr B}\).  Locality of \(\Phi\)
puts its image in \(\mathfrak m_B\), proving the
first inclusion in (G97.10).  Extension to \(A\)
gives the second.  The third is the defining
\(x_0\)-annihilation of every tuple \(\mathrm{FPt}\).
The polynomial families from AC-G95 are such maps
with \(Q=K[T]\), so the same proof applies to the
whole polynomial family, not merely its endpoints.
This is a statement about the closed invariant
fiber.  A map \(B\to K[[T]]\) with varying
weight-zero characters need not kill
\(\mathfrak m_B\); none of these inclusions proves
that it kills \(\Phi(J_{\rm aut})\).
\end{proof}

\paragraph{Source and geometric scope.}
The literal \path{NS}, \path{GS} and \path{toGS} are
in \path{ReciprocityTheory/GaloisValuations/Global.lean}.
The complete files \path{ContinuousResP.lean},
\path{UniversalOrdinaryDeterminantRing.lean}, and
\path{RepAlg.lean} were read for the native quotient,
\path{Resid}, \path{mu}, \path{Ip}, \path{Pair},
\path{GDet.IsCont}, polynomial-law and universal
map definitions.  The complete containing
\path{PropagationTheoremGL2.lean} was read for the
original \path{TwLink} restrictions.  All are at
the pinned source commit recorded in the manifest.
The standard local reciprocity and local-unit
structure inputs are retained explicitly; no new
proof of local class field theory or Lean build is
claimed.

Geometrically (G97.8) is the actual morphism from
the framed tuple chart to the native ordinary
determinant deformation space.  It preserves the
specified local characters and actual finite
inertia, with the source's cyclotomic polarization
normalization.  Equation (G97.10) places the closed
invariant fiber in the pulled-back proper support.
It does not establish a generic component or
varying-character curve in that support, the
numerical \path{Frame.Big} inequalities, or the
original counting budget.  This uses the existing
formal-deformation and finite-cover dictionary;
no new archival SGA reading is attributed.
